% Auto-generated Code Vault for 'BlackHoleDocument (Article)' [Latex]

\documentclass{article}
\usepackage[margin=1in]{geometry}
\usepackage{amsmath}
\usepackage{amsfonts}
\usepackage{amssymb}
\usepackage{graphicx}
\usepackage{float}
\usepackage{caption}
\usepackage{subcaption}

\title{The Mathematics of Black Holes and the Schwarzschild Radius}
\author{Anonymous}
\date{\today}

\begin{document}

\begin{titlepage}
    \centering
    \vspace*{2cm}
    {\huge \bf The Mathematics of Black Holes and the Schwarzschild Radius}
    \vspace{1cm}
    {\large Anonymous}
    \vspace{2cm}
    {\large \today}
\end{titlepage}

\begin{abstract}
This document provides an overview of the mathematics behind black holes, focusing on the Schwarzschild radius. The Schwarzschild radius is the radius of a sphere such that, if all the mass of an object were to be compressed within that sphere, the escape velocity from the surface of the sphere would equal the speed of light. The mathematical framework for understanding black holes is rooted in Einstein's theory of general relativity.
\end{abstract}

\section{Introduction to Black Holes}
Black holes are among the most fascinating objects in the universe, with properties that challenge our understanding of space and time. The point of no return, called the event horizon, is where the gravitational pull is so strong that not even light can escape. The event horizon is directly related to the Schwarzschild radius, which is given by the equation:
\[ r_s = \frac{2GM}{c^2} \]
where $r_s$ is the Schwarzschild radius, $G$ is the gravitational constant, $M$ is the mass of the object, and $c$ is the speed of light.

\section{Mathematical Framework}
The mathematical framework for understanding black holes is based on Einstein's theory of general relativity. The core equation of general relativity is the Einstein field equation, given by:
\[ R_{\mu\nu} - \frac{1}{2}Rg_{\mu\nu} = \frac{8\pi G}{c^4}T_{\mu\nu} \]
where $R_{\mu\nu}$ is the Ricci tensor, $R$ is the Ricci scalar, $g_{\mu\nu}$ is the metric tensor, $G$ is the gravitational constant, $c$ is the speed of light, and $T_{\mu\nu}$ is the stress-energy tensor.

For a spherically symmetric, non-rotating mass, the metric can be simplified to the Schwarzschild metric:
\[ ds^2 = \left(1 - \frac{2GM}{c^2r}\right)dt^2 - \frac{1}{c^2}\left(1 - \frac{2GM}{c^2r}\right)^{-1}dr^2 - r^2(d\theta^2 + \sin^2\theta d\phi^2) \]
where $ds^2$ is the interval element, $G$ is the gravitational constant, $M$ is the mass of the object, $c$ is the speed of light, and $r$ is the radial distance from the center of the object.

\section{Conclusion}
In conclusion, the mathematics of black holes, particularly the concept of the Schwarzschild radius, provides a deep insight into the nature of these cosmic phenomena. Through the lens of general relativity, we can understand the extreme gravitational environments that characterize black holes.

\begin{thebibliography}{9}
\bibitem{Einstein} Einstein, A. (1915). Die Grundlage der allgemeinen Relativitätstheorie. \textit{Annalen der Physik}, 49, 769-822.
\bibitem{Schwarzschild} Schwarzschild, K. (1916). On the Gravitational Field of a Mass Point According to Einstein's Theory. \textit{Sitzungsberichte der Königlich Preußischen Akademie der Wissenschaften (Berlin)}, 7, 189-196.
\end{thebibliography}

\end{document}
